\chapter{Introducción}

DePaso es una plataforma digital de logística urbana colaborativa orientada a PyMEs, comercios y clientes particulares del Área Metropolitana de Buenos Aires (AMBA). Su propósito es facilitar la gestión integral de envíos urbanos mediante una solución que permita solicitar, asignar, monitorear y completar entregas de manera eficiente, incorporando además herramientas para estimar el impacto ambiental asociado a cada operación.

A diferencia de las plataformas de envíos bajo demanda actualmente disponibles, que se basan principalmente en transportistas dedicados, DePaso propone un modelo híbrido que combina transporte exclusivo con transporte colaborativo. Este último aprovecha trayectos que los transportistas ya realizan por motivos propios, permitiendo utilizar capacidad disponible en desplazamientos existentes y reduciendo viajes adicionales mediante mecanismos de asignación que minimizan los desvíos.

La propuesta surge ante una serie de problemáticas presentes en la logística urbana actual:

\begin{itemize}[leftmargin=2em]
    \item Costos logísticos elevados para PyMEs y comercios que no cuentan con flota propia.
    \item Existencia de vehículos que realizan trayectos con capacidad disponible sin ser aprovechada.
    \item Ineficiencias en la distribución de última milla dentro de entornos urbanos densamente poblados.
    \item Emisiones de CO\textsubscript{2} asociadas a viajes dedicados que podrían evitarse mediante esquemas colaborativos.
    \item Escasas oportunidades para que personas que realizan desplazamientos cotidianos obtengan ingresos adicionales aprovechando sus recorridos habituales.
\end{itemize}

{\color{nuevo}La magnitud de estas problemáticas se refleja en datos del sector. La entrega de última milla concentra una fracción desproporcionada del costo logístico total y constituye el tramo más ineficiente de la cadena de distribución \parencite{YangEtAl2022}. En Argentina esa presión se intensifica: el Índice de Costos Logísticos elaborado por la Cámara Empresaria de Operadores Logísticos, homologado por la Universidad Tecnológica Nacional, acumuló un incremento del 22,2~\% en el primer semestre de 2026, muy por encima del 14~\% registrado en el mismo período del año anterior \parencite{CEDOL2025}. Esta ineficiencia tiende a agravarse: de mantenerse las tendencias actuales, se proyecta que para 2030 la demanda de entregas urbanas crezca un 78~\%, con un aumento del 36~\% en la cantidad de vehículos de reparto, del 32~\% en las emisiones de CO\textsubscript{2} y del 21~\% en la congestión de las principales ciudades \parencite{WEF2020}. En Argentina, el transporte explica cerca del 14~\% de las emisiones totales de gases de efecto invernadero del país, siendo la segunda fuente de emisiones dentro del sector energético \parencite{InventarioGEI2024}. En el plano local, un relevamiento de la Organización Internacional del Trabajo sobre trabajadores de plataformas de reparto en la Argentina encontró que el 95~\% de quienes habían ingresado recientemente lo hizo ante la dificultad de conseguir otro empleo; hacia fines de 2018, solo tres de las principales aplicaciones de reparto (Glovo, PedidosYa y Rappi) sumaban en conjunto más de 7.100 repartidores independientes registrados, la amplia mayoría bajo la figura de monotributista y sin relación de dependencia \parencite{OIT2020}.}

Frente a este escenario, DePaso busca articular una solución que aproveche la capacidad ociosa existente en los desplazamientos urbanos, al tiempo que ofrece alternativas de envío dedicado para los casos que así lo requieran. El proyecto se circunscribe a dicha área metropolitana y tiene como horizonte temporal el desarrollo de un prototipo funcional para finales de 2026. Esta combinación se traduce en un conjunto de modalidades de servicio que estructuran la operación de la plataforma y que se describen a continuación.

% ─────────────────────────────────────────────────────────────────────────────
\section{Modalidades de envío}

DePaso organiza sus servicios a partir de dos dimensiones: la forma en que se inicia la operación y el tipo de trayecto realizado por el transportista. La combinación de ambas da lugar a cuatro modalidades de servicio:

\begin{itemize}[leftmargin=2em]
    \item \textbf{Dedicado por demanda}: el remitente solicita un envío y el sistema identifica transportistas disponibles cercanos al punto de origen.
    
    \item \textbf{Dedicado por espacio}: el transportista informa disponibilidad en una zona y franja horaria determinadas, permitiendo la asignación de pedidos compatibles con la capacidad de su vehículo.
    
    \item \textbf{Colaborativo por demanda}: el remitente solicita un envío con cierta flexibilidad temporal y el sistema busca trayectos existentes compatibles con el recorrido requerido, procurando minimizar el desvío adicional.
    
    \item \textbf{Colaborativo por espacio}: el transportista registra trayectos habituales recurrentes y recibe sugerencias de envíos compatibles con dichos recorridos.
\end{itemize}

Para la operación del sistema se contemplan cuatro categorías generales de carga: paquetes pequeños y documentos, cargas medianas, cargas grandes o voluminosas y mudanzas o fletes. Las restricciones según el tipo de movilidad disponible se presentan en la Tabla~\ref{tab:restricciones}.
\begin{table}[H]
    \caption{Restricciones de carga por tipo de movilidad}
    \label{tab:restricciones}
    \centering
    \small
    \begin{tabular}{lccccc}
        \toprule
        \textbf{Categoría} & \textbf{Peatón/Bici} & \textbf{Moto} & \textbf{Auto} & \textbf{Camioneta} & \textbf{Camión} \\
        \midrule
        Pequeños / documentos & Sí\textsuperscript{*} & Sí & Sí & Sí & Sí \\
        Medianos              & No & Sí & Sí & Sí & Sí \\
        Grandes / voluminosos & No & No & Sí & Sí & Sí \\
        Mudanzas / fletes     & No & No & No & Sí & Sí \\
        \bottomrule
        \multicolumn{6}{l}{\textsuperscript{*} Solo trayectos cortos ($<$ 5 km).}
    \end{tabular}
\end{table}

Las mudanzas y fletes se asignan siempre en modalidad dedicada; nunca colaborativa.

% ─────────────────────────────────────────────────────────────────────────────
\vspace{1em}
\section{Objetivos}

\subsection{Objetivo general}

Desarrollar, para fines de 2026, un prototipo funcional de plataforma digital de logística urbana colaborativa para el AMBA que permita gestionar el ciclo completo de los envíos mediante mecanismos inteligentes de clasificación y asignación, aprovechando la capacidad disponible en los desplazamientos de los transportistas.

\subsection{Objetivos específicos}

{\color{nuevo}
\begin{itemize}[leftmargin=2em]
    \item Desarrollar e implementar un sistema de visión computacional capaz de clasificar automáticamente las cargas en las cuatro categorías volumétricas definidas a partir de imágenes, alcanzando una exactitud (\emph{accuracy}) sobre el conjunto de prueba de al menos el 60\%, valor que supera ampliamente la línea base aleatoria del 25\% correspondiente a cuatro categorías.

    \item Implementar un mecanismo de asignación inteligente que evalúe la compatibilidad entre pedidos y transportistas considerando características de la carga, trayectorias, ventanas horarias y reputación, y que restrinja el desvío del transportista colaborativo a un máximo del 15\% del trayecto original mediante un filtro duro, devolviendo un ranking de candidatos con su puntaje desglosado y explicable.

    \item Incorporar un modelo híbrido de operación que contemple tanto envíos dedicados como colaborativos, cubriendo las cuatro categorías de carga y los distintos tipos de movilidad definidos.

    \item Desarrollar un mecanismo de estimación del impacto ambiental que calcule el ahorro de CO\textsubscript{2} del 100\% de los envíos colaborativos completados respecto de un escenario equivalente de transporte dedicado, a partir de factores de emisión estandarizados por tipo de vehículo, verificando el cálculo mediante pruebas automatizadas con coordenadas reales del AMBA.

    \item Diseñar e implementar una plataforma digital con una arquitectura modular preparada para escalar, que integre los distintos componentes funcionales necesarios para la gestión de envíos urbanos, con una cobertura geográfica objetivo del AMBA y con pruebas automatizadas sobre la totalidad de los módulos críticos, incluidos los flujos de integración de punta a punta.
\end{itemize}
}

% ─────────────────────────────────────────────────────────────────────────────
\vspace{1em}
\section{Alcance}

El proyecto comprende el diseño y desarrollo de un prototipo funcional de plataforma digital para la gestión logística de envíos dentro del AMBA.

\subsection{Características del prototipo funcional}

El sistema incluirá:

\begin{itemize}[leftmargin=2em]
    \item Registro y autenticación de clientes y transportistas con distintos medios de movilidad.
    
    \item Creación, asignación y seguimiento de pedidos en las diferentes modalidades de servicio definidas.
    
    \item Clasificación automática de cargas mediante inteligencia artificial, con posibilidad de ingreso manual alternativo.
    
    \item Asignación inteligente basada en trayectorias geográficas compatibles.
    
    \item Gestión de capacidad disponible de los vehículos.
    
    \item Estimación del ahorro de CO\textsubscript{2} asociado a envíos colaborativos.
    
    \item Herramientas básicas de monitoreo operativo para administradores.
\end{itemize}

\vspace{1em}
\subsection{Limitaciones y exclusiones}

El alcance se restringe geográficamente a operaciones dentro del AMBA, excluyendo envíos interprovinciales e internacionales. El clasificador de carga estará acotado a las categorías volumétricas definidas y al dataset de entrenamiento, ofreciendo el ingreso manual como alternativa cuando el objeto no sea reconocido con confianza suficiente. La pasarela de pagos se simulará en el prototipo, dejando su integración productiva como trabajo posterior. El mecanismo de cobertura ante daños se abordará a nivel de diseño y definición de responsabilidades, quedando la integración con compañías aseguradoras fuera del alcance del prototipo.

\subsection{Estructura del documento}

El presente documento se organiza de la siguiente manera: el Capítulo~1 introduce el problema, las modalidades de servicio, los objetivos y el alcance del proyecto; el Capítulo~2 presenta los antecedentes, compuestos por el marco teórico, el estado del arte y {\color{nuevo}el análisis competitivo}; el Capítulo~3 describe la propuesta y la investigación de usuarios realizada; {\color{nuevo}el Capítulo~4 detalla el análisis de requerimientos, con los requerimientos funcionales y no funcionales y los casos de uso; el Capítulo~5 desarrolla la metodología de desarrollo, la arquitectura y las tecnologías, el modelo de datos y la validación del sistema; y el Capítulo~6} expone la conclusión parcial y el trabajo futuro. Completan el documento la bibliografía y los anexos, que incluyen el cronograma de actividades{\color{nuevo}, la transcripción de los resultados de la encuesta y la transcripción de las entrevistas semiestructuradas}.